\documentclass[11pt,letterpaper]{article}

% ==============================================================================
% PAQUETES FUNDAMENTALES Y CONFIGURACIÓN REGIONAL
% ==============================================================================
\usepackage[utf8]{inputenc}
\usepackage[T1]{fontenc}
\usepackage[spanish,es-noshorthands]{babel}
\usepackage[margin=2.1cm, top=2.3cm, bottom=2.3cm, headheight=15pt]{geometry}
\usepackage{amsmath,amssymb,amsfonts,bm}
\usepackage{graphicx}
\usepackage{booktabs}
\usepackage{tabularx}
\usepackage{enumitem}
\usepackage{microtype}
\usepackage{xcolor}
\usepackage{fancyhdr}
\usepackage{lastpage}
\usepackage{float}

% ==============================================================================
% DEFINICIÓN DE PALETA DE COLOR INSTITUCIONAL UIS Y ACENTO
% ==============================================================================
\definecolor{uisdark}{HTML}{1A1D20}
\definecolor{uispurple}{HTML}{3D2B56}
\definecolor{uisgreen}{HTML}{1F6F43}
\definecolor{uisblue}{HTML}{0D47A1}
\definecolor{uisamber}{HTML}{B76E00}
\definecolor{uisrose}{HTML}{881337}
\definecolor{uislightbg}{HTML}{F8FAFC}
\definecolor{uisborder}{HTML}{E2E8F0}
\definecolor{uisgray}{HTML}{64748B}

% ==============================================================================
% TIKZ Y DIAGRAMAS VECTORIALES
% ==============================================================================
\usepackage{tikz}
\usetikzlibrary{arrows.meta, patterns, calc, positioning, decorations.pathmorphing, shapes.geometric, babel}

% ==============================================================================
% HIPERVÍNCULOS Y METADATOS
% ==============================================================================
\usepackage{hyperref}
\hypersetup{
    colorlinks=true,
    linkcolor=uispurple,
    citecolor=uisblue,
    urlcolor=uispurple,
    pdftitle={Programa y Contenidos de Asignatura: Análisis No Lineal de Estructuras},
    pdfauthor={Prof. Orlando Arroyo, Ph.D. - UIS},
    pdfsubject={Programa Curricular y Guía Temática de Posgrados},
    pdfkeywords={Análisis No Lineal, OpenSeesPy, Moodle, Pushover, Método N2, Fibras, MVLEM}
}

% ==============================================================================
% ENCABEZADOS Y PIES DE PÁGINA (FANCYHDR)
% ==============================================================================
\pagestyle{fancy}
\fancyhf{}
\fancyhead[L]{\small\color{uisgray}\textsf{Análisis No Lineal de Estructuras}}
\fancyhead[R]{\small\color{uisgray}\textsf{Programa de Asignatura}}
\fancyfoot[L]{\small\color{uisgray}\textsf{Escuela de Ingeniería Civil}}
\fancyfoot[C]{\small\color{uisgray}\textsf{Posgrados}}
\fancyfoot[R]{\small\color{uisgray}\textsf{Pág. \thepage\ de \pageref{LastPage}}}
\renewcommand{\headrulewidth}{0.4pt}
\renewcommand{\footrulewidth}{0.4pt}

% ==============================================================================
% COMIENZO DEL DOCUMENTO
% ==============================================================================
\begin{document}

% --- ENCABEZADO TITULAR INSTITUCIONAL ---
\begin{center}
    {\scshape\Large Universidad Industrial de Santander}\\[0.18cm]
    {\large Escuela de Ingeniería Civil}\\[0.12cm]
    {\normalsize Posgrados}\\[0.3cm]
    {\color{uispurple}\hrule height 2pt}\vspace{0.35cm}
    
    {\Huge\bfseries\color{uisdark} Análisis No Lineal de Estructuras}\\[0.2cm]
    {\Large\bfseries\color{uispurple} Programa de Asignatura y Contenidos Temáticos}\\[0.25cm]
    
    {\color{uispurple}\hrule height 1pt}\vspace{0.3cm}
    
    \small
    \begin{tabularx}{\textwidth}{@{} l X l X @{}}
        \textbf{Profesor:} & Prof. Orlando Arroyo, Ph.D. & \textbf{Periodo:} & Semestre II \\
        \textbf{Créditos:} & 4 Créditos (Posgrados) & \textbf{Modalidad:} & Teórico-Práctica Computacional \\
        \textbf{Software:} & OpenSeesPy (\texttt{openseespy}) & \textbf{Plataforma:} & Moodle UIS
    \end{tabularx}
\end{center}

\vspace{0.15cm}

% ==============================================================================
% SECCIÓN 1: INTRODUCCIÓN Y RECURSOS DEL CURSO
% ==============================================================================
\section{Presentación del Curso, Plataformas y Recursos Digitales}

El curso de \textbf{Análisis No Lineal de Estructuras} constituye el núcleo avanzado de formación en ingeniería sísmica y resiliencia estructural para los programas de Posgrados de la Escuela de Ingeniería Civil de la Universidad Industrial de Santander (UIS). La asignatura trasciende el paradigma convencional de diseño basado en fuerzas y coeficientes elásticos globales de reducción ($R$), dotando al estudiante de los fundamentos matemáticos, mecánicos y computacionales requeridos para simular el comportamiento inelástico real de edificaciones de concreto reforzado bajo acciones dinámicas extremas.

Para asegurar una experiencia académica moderna, interactiva y reproducible, la asignatura se apoya en las siguientes plataformas y recursos oficiales:

\begin{enumerate}[leftmargin=16pt, topsep=3pt, itemsep=3pt, parsep=1pt, label=\textbf{\arabic*.}]
    \item \textbf{Plataforma Web Interactiva del Curso (Landing Page):}
    \begin{itemize}[leftmargin=12pt, topsep=1pt, itemsep=1pt]
        \item \textbf{URL Oficial:} \href{https://odarroyo.github.io/analisis_no_lineal/}{\texttt{https://odarroyo.github.io/analisis\_no\_lineal/}}
        \item \textit{Propósito:} Aloja los laboratorios numéricos interactivos de acceso libre desarrollados específicamente para el curso. Permite a los estudiantes interactuar directamente con simuladores cinemáticos en tiempo real (visores de dependencia de trayectoria, analizadores del residuo de Newton-Raphson, inspección matricial paso a paso y gráficas interactivas).
    \end{itemize}

    \item \textbf{Curso Complementario de Concreto Reforzado:}
    \begin{itemize}[leftmargin=12pt, topsep=1pt, itemsep=1pt]
        \item \textbf{URL de Consulta:} \href{https://odarroyo.github.io/curso_concreto_reforzado/}{\texttt{https://odarroyo.github.io/curso\_concreto\_reforzado/}}
        \item \textit{Propósito:} Repositorio complementario sobre fundamentos de diseño sísmico de secciones, diagramas de interacción $P-M$, confinamiento seccional y detallamiento dúctil, necesarios para la adecuada calibración de modelos inelásticos.
    \end{itemize}

    \item \textbf{Plataforma Computacional Oficial --- OpenSeesPy:}
    \begin{itemize}[leftmargin=12pt, topsep=1pt, itemsep=1pt]
        \item \textit{Entorno de Trabajo:} Toda la componente numérica y los proyectos de análisis no lineal se desarrollarán mediante la librería en Python de \textbf{OpenSees} (\textit{Open System for Earthquake Engineering Simulation}): \texttt{openseespy}.
        \item \textit{Integración Científica:} Facilita el uso del ecosistema científico de Python (\texttt{NumPy}, \texttt{SciPy}, \texttt{Matplotlib}, \texttt{opsvis}, \texttt{vfo}) para preprocesamiento automatizado, solución paramétrica y visualización de resultados.
    \end{itemize}

    \item \textbf{Plataforma Virtual Institucional (Moodle UIS):}
    \begin{itemize}[leftmargin=12pt, topsep=1pt, itemsep=1pt]
        \item Los estudiantes matriculados formalmente en la asignatura serán incorporados al espacio virtual en \textbf{Moodle UIS}. En esta plataforma se centralizarán las grabaciones de clase, enunciados y entregas del proyecto de curso, rúbricas de evaluación, foros técnicos y literatura de consulta.
    \end{itemize}
\end{enumerate}

\vspace{0.2cm}

% ==============================================================================
% SECCIÓN 2: OBJETIVOS DE APRENDIZAJE
% ==============================================================================
\section{Objetivos Formativos}

\subsection{Objetivo General}
Capacitar al estudiante en la formulación rigurosa, implementación numérica, análisis computacional y juicio crítico del comportamiento inelástico y no lineal de estructuras de concreto reforzado sometidas a demandas sísmicas severas, dominando las herramientas computacionales del estado del arte (OpenSeesPy) y las metodologías de evaluación basadas en desempeño sismorresistente.

\subsection{Objetivos Específicos}
\begin{itemize}[leftmargin=16pt, topsep=3pt, itemsep=2.5pt, parsep=1pt]
    \item \textbf{OE1 (Fundamentos Numéricos):} Comprender los principios matemáticos que gobiernan el equilibrio no lineal, la no unicidad de soluciones, la dependencia de trayectoria y la aplicación rigurosa de algoritmos de resolución incremental-iterativa (Newton-Raphson estándar y modificado, control por desplazamiento y longitud de arco).
    \item \textbf{OE2 (Fuentes de No Linealidad):} Distinguir, formular y evaluar las contribuciones de la no linealidad material (fluencia, degradación, agrietamiento) y geométrica ($P-\Delta$, $P-\delta$ y transformaciones corotacionales).
    \item \textbf{OE3 (Modelación de Materiales):} Calibrar leyes constitutivas uniaxiales y multiaxiales avanzadas para concreto confinado y no confinado (Mander) y acero de refuerzo con efecto Bauschinger (Menegotto-Pinto / Giuffré-Pinto).
    \item \textbf{OE4 (Pórticos de Concreto):} Formular e implementar modelos de elementos finitos de plasticidad concentrada (resortes $M-\theta$ con modelos histeréticos de daño tipo Ibarra-Medina-Krawinkler) y de plasticidad distribuida en fibras (elementos basados en desplazamientos y basados en fuerzas).
    \item \textbf{OE5 (Muros Estructurales):} Formular y analizar sistemas de muros mediante elementos laminares multicapa (\textit{Multilayer Shell}), modelos de fibras y macroelementos con interacción acoplada flexión-corte (MVLEM y SFI-MVLEM).
    \item \textbf{OE6 (Evaluación por Desempeño):} Ejecutar análisis estáticos no lineales pushover, transformar sistemas continuos a osciladores equivalentes SDOF e implementar el \textbf{Método N2 Modificado} en formato espectral ADRS para determinar el punto de desempeño sísmico.
\end{itemize}

\vspace{0.2cm}

% --- MAPA JERÁRQUICO DEL CURSO (TIKZ) ---
\begin{figure}[H]
\centering
\resizebox{0.88\linewidth}{!}{%
\begin{tikzpicture}[
    node distance=0.95cm and 0.35cm,
    block/.style={rectangle, draw=uisdark, fill=uislightbg, rounded corners=3pt, line width=1.1pt, text width=3.2cm, align=center, font=\small\sffamily, minimum height=1.0cm},
    headerblock/.style={rectangle, draw=uispurple, fill=uispurple!10, rounded corners=3pt, line width=1.3pt, text width=3.3cm, align=center, font=\small\sffamily\bfseries, minimum height=1.0cm},
    mainnode/.style={rectangle, draw=uisblue, fill=uisblue!15, rounded corners=4pt, line width=1.4pt, text width=11.2cm, align=center, font=\normalsize\sffamily\bfseries, minimum height=0.8cm},
    arr/.style={->, >=Stealth, thick, uispurple}
]

\node[mainnode] (fund) {Fundamento Central: Equilibrio Inelástico Dependiente de la Trayectoria y OpenSeesPy};

\node[headerblock, below left=0.95cm and -0.2cm of fund] (mats) {\textbf{No Linealidad de Material}\\[2pt]\footnotesize Modelos Concreto / Acero};
\node[headerblock, below=0.95cm of fund] (elems) {\textbf{Formulación de Elementos}\\[2pt]\footnotesize Fibras, Rótulas, MVLEM};
\node[headerblock, below right=0.95cm and -0.2cm of fund] (geom) {\textbf{No Linealidad Geométrica}\\[2pt]\footnotesize Efectos $P-\Delta$, $P-\delta$, Corotacional};

\draw[arr] (fund.south) -- ++(0,-0.3) -| (mats.north);
\draw[arr] (fund.south) -- (elems.north);
\draw[arr] (fund.south) -- ++(0,-0.3) -| (geom.north);

\node[block, below=0.8cm of elems, fill=uisamber!15, draw=uisamber, line width=1.3pt, text width=9.2cm] (eval) {\textbf{\color{uisdark}Evaluación Sísmica No Lineal y Desempeño Estructural}\\[2pt]\small Análisis Pushover Multi-modal y Método N2 Modificado (ADRS / Eurocódigo 8 / ASCE 41)};

\draw[arr, uisamber] (mats.south) |- (eval.west);
\draw[arr, uisamber] (elems.south) -- (eval.north);
\draw[arr, uisamber] (geom.south) |- (eval.east);

\end{tikzpicture}%
}
\caption{Estructura conceptual y articulación temática de la asignatura.}
\label{fig:course_map}
\end{figure}

\vspace{0.2cm}

% ==============================================================================
% SECCIÓN 3: CONTENIDOS TEMÁTICOS (LISTA DE LO QUE HAREMOS)
% ==============================================================================
\section{Contenidos Temáticos de la Asignatura}

El curso se desarrollará a través de los siguientes 7 módulos temáticos, estructurados como una guía directa de las actividades analíticas, teóricas y computacionales que realizaremos:

\begin{enumerate}[leftmargin=16pt, topsep=2pt, itemsep=3.5pt, parsep=1pt, label=\textbf{\arabic*.}]

    % --- TEMA 1 ---
    \item \textbf{¿Qué es el análisis no lineal?}
    \begin{itemize}[leftmargin=12pt, topsep=1pt, itemsep=1.5pt]
        \item Comprender por qué la trayectoria de carga importa en estructuras que incursionan en el rango inelástico (\textit{path-dependence}).
        \item Demostrar la pérdida de validez del principio de superposición y la no conmutatividad de combinaciones de carga mediante osciladores elastoplásticos con endurecimiento.
        \item Plantear la ecuación fundamental de equilibrio en términos del residuo de fuerzas: $\mathbf{r}(\mathbf{u}) = \mathbf{P} - \mathbf{F}_{\text{int}}(\mathbf{u}) = \mathbf{0}$.
        \item Implementar el algoritmo iterativo de Newton-Raphson estándar con cálculo y actualización explícita de la matriz de rigidez tangente $\mathbf{K}_t$.
        \item Analizar variantes algorítmicas de resolución: Newton-Raphson modificado (rigidez constante $\mathbf{K}_0$) y esquemas Quasi-Newton (BFGS, Broyden, Krylov-Newton).
        \item Evaluar normas y criterios de convergencia: norma de desequilibrio residual ($\|\mathbf{r}\|_2$), norma de incrementos de desplazamiento ($\|\delta\mathbf{u}\|_2$) y norma de energía interna ($\frac{1}{2}\delta\mathbf{u}^T\mathbf{r}$).
        \item Aplicar estrategias de solución ante inestabilidades, reblandecimiento numérico y puntos límite: control por desplazamiento (\textit{displacement control}) y métodos de longitud de arco (\textit{arc-length / Riks-Crisfield}).
    \end{itemize}

    % --- TEMA 2 ---
    \item \textbf{Tipos de análisis estructural: Lineal vs. No lineal, Estático vs. Dinámico}
    \begin{itemize}[leftmargin=12pt, topsep=1pt, itemsep=1.5pt]
        \item Estructurar la matriz taxonómica $2\times 2$ de los procedimientos de análisis estructural:
        \begin{itemize}[leftmargin=10pt, topsep=1pt, itemsep=1pt]
            \item \textit{Lineal Estático:} Método de fuerzas horizontales equivalentes (ELF).
            \item \textit{Lineal Dinámico:} Análisis modal espectral y análisis en el tiempo elástico (\textit{Elastic RHA}).
            \item \textit{No Lineal Estático:} Análisis pushover con patrones laterales de carga.
            \item \textit{No Lineal Dinámico:} Análisis cronológico paso a paso (\textit{Nonlinear Time History Analysis} - NL-THA).
        \end{itemize}
        \item Contrastar los rangos de validez, hipótesis cinemáticas, costos computacionales y limitaciones de cada procedimiento.
        \item Analizar el significado físico y las limitaciones de los factores de reducción de fuerza por ductilidad ($R$) y sobrerresistencia ($\Omega_0$) en el diseño convencional.
        \item Revisar las prescripciones normativas para la aplicabilidad del análisis no lineal (NSR-10, ASCE 7, Eurocódigo 8).
    \end{itemize}

    % --- TEMA 3 ---
    \item \textbf{Fuentes de no linealidad}
    \begin{enumerate}[leftmargin=14pt, topsep=1pt, itemsep=2pt, label=\textbf{3.\arabic*.}]
        \item \textbf{No linealidad de material:}
        \begin{itemize}[leftmargin=10pt, topsep=1pt, itemsep=1pt]
            \item Estudiar el agrietamiento por tracción del concreto y la pérdida abrupta de rigidez elástica seccional.
            \item Modelar la fluencia del acero de refuerzo longitudinal y el endurecimiento por deformación (\textit{strain hardening}).
            \item Simular el aplastamiento a compresión y la degradación post-pico del concreto (\textit{strain softening}).
            \item Analizar la pérdida de adherencia (\textit{bond-slip}), degradación de rigidez/resistencia y estrangulamiento histerético (\textit{pinching}).
        \end{itemize}
        \item \textbf{No linealidad geométrica:}
        \begin{itemize}[leftmargin=10pt, topsep=1pt, itemsep=1pt]
            \item Formular y evaluar el efecto de segundo orden global $P-\Delta$ (amplificación de momentos de entrepiso).
            \item Formular el efecto de segundo orden local de miembro $P-\delta$ en columnas esbeltas sometidas a flexocompresión.
            \item Deducir y ensamblar la matriz de rigidez geométrica $\mathbf{K}_G$ y evaluar condiciones de inestabilidad / pandeo.
            \item Implementar transformaciones geométricas en OpenSeesPy: lineal (\texttt{Linear}), $P-\Delta$ (\texttt{PDelta}) y corotacional (\texttt{Corotational}).
        \end{itemize}
        \item \textbf{Modelos no lineales de materiales:}
        \begin{itemize}[leftmargin=10pt, topsep=1pt, itemsep=1pt]
            \item Calibrar el modelo de Mander et al. para concreto confinado y no confinado (curvas monotónicas y cíclicas).
            \item Implementar materiales en OpenSeesPy: \texttt{Concrete01} (sin tracción), \texttt{Concrete02} (con degradación a tracción) y \texttt{Concrete04}.
            \item Implementar el modelo de Menegotto-Pinto / Giuffré-Pinto para acero de refuerzo con efecto Bauschinger (\texttt{Steel02}).
            \item Realizar pruebas computacionales de respuesta constitutiva bajo historias cíclicas de deformación impuesta.
        \end{itemize}
    \end{enumerate}

    % --- TEMA 4 ---
    \item \textbf{Formulaciones no lineales para elementos de pórticos de concreto reforzado}
    \begin{enumerate}[leftmargin=14pt, topsep=1pt, itemsep=2pt, label=\textbf{4.\arabic*.}]
        \item \textbf{Modelos de plasticidad concentrada:}
        \begin{itemize}[leftmargin=10pt, topsep=1pt, itemsep=1pt]
            \item Modelar vigas y columnas mediante elementos elásticos centrales con resortes rotacionales de longitud nula (\textit{zeroLength}) en los extremos.
            \item Construir curvas analíticas momento-curvatura ($M-\phi$) e integrarlas en relaciones momento-rotación ($M-\theta$).
            \item Estimar la longitud de rótula plástica equivalente ($L_p$) según formulaciones empíricas (Park-Paulay, Priestley).
            \item Implementar modelos histeréticos con degradación cíclica tipo Ibarra-Medina-Krawinkler (IMK / \texttt{ModIMKPinching}).
        \end{itemize}
        \item \textbf{Modelos de fibras:}
        \begin{itemize}[leftmargin=10pt, topsep=1pt, itemsep=1pt]
            \item Discretizar secciones transversales en mallas de fibras uniaxiales (núcleo confinado, recubrimiento y barras de acero).
            \item Implementar elementos basados en desplazamientos (\texttt{dispBeamColumn}) y evaluar su sensibilidad al mallado.
            \item Implementar elementos formulados en fuerzas (\texttt{forceBeamColumn}) con integración numérica de Gauss-Lobatto.
            \item Tratar la localización de deformaciones y regularización de longitud de rótula plástica en ablandamiento.
        \end{itemize}
    \end{enumerate}

    % --- TEMA 5 ---
    \item \textbf{Formulaciones para elementos de edificios de muros}
    \begin{enumerate}[leftmargin=14pt, topsep=1pt, itemsep=2pt, label=\textbf{5.\arabic*.}]
        \item \textbf{Formulaciones Shell Multilayer:}
        \begin{itemize}[leftmargin=10pt, topsep=1pt, itemsep=1pt]
            \item Discretizar muros estructurales mediante elementos finitos laminares multicapa 2D (\texttt{ShellMITC4}, \texttt{ShellDKGQ}).
            \item Definir capas ortótropas superpuestas de concreto y refuerzo longitudinal/transversal en estado plano de esfuerzos.
            \item Simular distorsión por cortante, agrietamiento diagonal y pandeo de borde mediante la teoría de campo de compresión modificada (MCFT).
        \end{itemize}
        \item \textbf{Formulación de fibras (columna ancha / beam-column):}
        \begin{itemize}[leftmargin=10pt, topsep=1pt, itemsep=1pt]
            \item Modelar muros mediante elementos viga-columna de fibras colocados en el eje neutroide conectados a diafragmas mediante brazos rígidos (\textit{rigid links}).
            \item Analizar la captura del acoplamiento carga axial-flexión y discutir el desacoplamiento del cortante mediante resortes externos.
        \end{itemize}
        \item \textbf{Formulaciones de MVLEM y SFI-MVLEM:}
        \begin{itemize}[leftmargin=10pt, topsep=1pt, itemsep=1pt]
            \item Implementar el macroelemento de múltiples líneas verticales (\textit{Multiple-Vertical-Line-Element Model} --- MVLEM) con macrofibras uniaxiales y resorte de corte horizontal desacoplado.
            \item Implementar el macroelemento con interacción acoplada corte-flexión (\textit{Shear-Flexure Interaction} MVLEM --- SFI-MVLEM) en OpenSeesPy.
            \item Analizar y predecir modos de falla: aplastamiento de talón por flexocompresión, fluencia por flexión, falla por corte y deslizamiento basal.
        \end{itemize}
    \end{enumerate}

    % --- TEMA 6 ---
    \item \textbf{Análisis pushover (estático no lineal)}
    \begin{itemize}[leftmargin=12pt, topsep=1pt, itemsep=1.5pt]
        \item Definir la secuencia metodológica de aplicación de cargas: gravitacionales de servicio ($1.0D + 0.25L$) bajo control de carga seguidas del empuje lateral incremental bajo control de desplazamiento.
        \item Configurar y comparar patrones de distribución lateral de fuerzas: uniforme por masa, triangular/modal basado en el primer modo elástico y patrones multimodales.
        \item Extraer y construir la curva de capacidad global de la edificación (cortante en la base $V_b$ vs. desplazamiento en la cubierta $\Delta_{\text{techo}}$).
        \item Transformar rigurosamente el sistema continuo (MDOF) al oscilador equivalente de un grado de libertad (SDOF) mediante el factor de participación modal $\Gamma$ y la masa modal efectiva $M^*$.
        \item Identificar la secuencia de plastificación progresiva, formación de articulaciones plásticas, sobrerresistencia seccional y mecanismo cinemático de colapso.
    \end{itemize}

    % --- TEMA 7 ---
    \item \textbf{Método N2 modificado}
    \begin{itemize}[leftmargin=12pt, topsep=1pt, itemsep=1.5pt]
        \item Revisar los fundamentos analíticos del Método N2 (Fajfar, 2000) y su adopción en estándares de evaluación sismorresistente (Eurocódigo 8).
        \item Transformar la curva de capacidad del oscilador SDOF y el espectro de respuesta elástico al formato espectral de aceleración-desplazamiento (ADRS: $S_a$ vs. $S_d$).
        \item Realizar la idealización bilineal elasto-plástica de la curva de capacidad mediante balance de áreas de energía de deformación plástica.
        \item Determinar el factor de reducción inelástica por ductilidad $R_\mu$ y clasificar la respuesta según el período elástico característico del sistema ($T^*$).
        \item Calcular el Punto de Desempeño ($d^*$) mediante la regla de igual desplazamiento ($T^* \ge T_C$) o la regla de igual energía ($T^* < T_C$).
        \item Aplicar la extensión modificada para considerar los efectos de torsión acoplada y la amplificación de demandas por modos superiores en la altura de la estructura.
        \item Evaluar y verificar los estados límite de daño y niveles de desempeño estructural (Operacional, Ocupación Inmediata, Seguridad de Vida, Prevención de Colapso) de conformidad con ASCE/SEI 41-17 y Eurocódigo 8.
    \end{itemize}

\end{enumerate}

\vspace{0.25cm}

% ==============================================================================
% SECCIÓN 4: METODOLOGÍA Y EVALUACIÓN
% ==============================================================================
\section{Metodología de Enseñanza y Esquema de Evaluación}

La metodología de la asignatura articula el desarrollo conceptual riguroso en las sesiones magistrales con talleres prácticos guiados de programación e implementación computacional en \textbf{OpenSeesPy}. La plataforma virtual de \textbf{Moodle UIS} constituirá el entorno oficial para la publicación de material docente, foros de soporte técnico y gestión de entregas académicas.

El esquema de evaluación se concentra en el desarrollo progresivo de un \textbf{Proyecto Práctico de Modelación No Lineal}, el cual cuenta con dos entregas principales que concentran el \textbf{90\%} de la calificación definitiva, complementado por un \textbf{Quiz Final} conceptual con valor del \textbf{10\%}.

\vspace{0.2cm}
\begin{table}[H]
\centering
\caption{Esquema oficial de evaluación de la asignatura.}
\label{tab:evaluacion}
\small
\setlength{\tabcolsep}{5pt}
\renewcommand{\arraystretch}{1.3}
\begin{tabularx}{\textwidth}{l p{8.2cm} c c}
\toprule
\textbf{Actividad Evaluativa} & \textbf{Descripción, Alcance y Entregables en Moodle UIS} & \textbf{Porcentaje} & \textbf{Semanas} \\
\midrule
\textbf{Proyecto: Entrega 1} & 
\textbf{Modelación Constitutiva Seccional y de Componentes en OpenSeesPy:} \newline
Calibración de leyes constitutivas no lineales de materiales (concreto confinado/no confinado según Mander y acero con efecto Bauschinger según Menegotto-Pinto). Modelación comparativa de elementos de pórticos (plasticidad concentrada vs. fibras) y muros de corte (macroelementos MVLEM / SFI-MVLEM) con verificación experimental. & 
\textbf{45\%} & Sem. 8--9 \\
\midrule
\textbf{Proyecto: Entrega 2} & 
\textbf{Modelo Global, Análisis Pushover y Método N2 Modificado:} \newline
Construcción y ensamblaje del modelo inelástico completo de una edificación de concreto reforzado. Ejecución del análisis pushover con patrones modales y uniformes. Conversión al oscilador SDOF, aplicación del Método N2 Modificado en formato ADRS, localización del Punto de Desempeño y evaluación de estados límite según ASCE 41-17 y Eurocódigo 8. & 
\textbf{45\%} & Sem. 15--16 \\
\midrule
\textbf{Quiz Final} & 
\textbf{Evaluación Conceptual de Fundamentos y Algoritmos:} \newline
Prueba individual sobre principios matemáticos del análisis no lineal, residuo de equilibrio, algoritmos de Newton-Raphson, matrices tangentes, criterios de convergencia, formulaciones de elementos finitos y metodologías de desempeño. & 
\textbf{10\%} & Sem. 16 \\
\midrule
\textbf{Total} & \textbf{Evaluación Integral de la Asignatura} & \textbf{100\%} & \\
\bottomrule
\end{tabularx}
\end{table}

\vspace{0.25cm}

% ==============================================================================
% SECCIÓN 5: REFERENCIAS BIBLIOGRÁFICAS
% ==============================================================================
\section{Bibliografía y Literatura de Referencia}

\begin{enumerate}[leftmargin=16pt, topsep=2pt, itemsep=1.5pt, label={[\arabic*]}]
    \item \textbf{Chopra, A. K. (2020).} \textit{Dynamics of Structures: Theory and Applications to Earthquake Engineering}. 5th Edition, Pearson Education.
    \item \textbf{Priestley, M. J. N., Calvi, G. M., \& Kowalsky, M. J. (2007).} \textit{Displacement-Based Seismic Design of Structures}. IUSS Press, Pavia, Italy.
    \item \textbf{Powell, G. H. (2010).} \textit{Nonlinear Analysis of Structures for Seismic Design: Outline of Theory and Computation}. Computers and Structures Inc., Berkeley, CA.
    \item \textbf{Fajfar, P. (2000).} ``A Nonlinear Analysis Method for Performance-Based Seismic Design''. \textit{Earthquake Spectra}, 16(3), 573--592.
    \item \textbf{Kolozvari, K., Orakcal, K., \& Wallace, J. W. (2015).} ``Modeling of Cyclic Shear-Flexure Interaction in Reinforced Concrete Structural Walls. I: Theory''. \textit{Journal of Structural Engineering}, 141(5), 04014135.
    \item \textbf{Mander, J. B., Priestley, M. J. N., \& Park, R. (1988).} ``Theoretical Stress-Strain Model for Confined Concrete''. \textit{Journal of Structural Engineering}, 114(8), 1804--1826.
    \item \textbf{Menegotto, M., \& Pinto, P. E. (1973).} ``Method of analysis for cyclically loaded RC frames including changes in geometry and non-elastic behavior of elements under combined normal force and bending''. \textit{IABSE Symposium on Resistance and Ultimate Deformability of Structures Acted on by Well-Defined Repeated Loads}, Lisbon, Portugal.
    \item \textbf{ASCE/SEI 41-17 (2017).} \textit{Seismic Evaluation and Retrofit of Existing Buildings}. American Society of Civil Engineers, Reston, Virginia.
    \item \textbf{European Committee for Standardization (CEN) (2004).} \textit{Eurocode 8: Design of structures for earthquake resistance - Part 1: General rules, seismic actions and rules for buildings (EN 1998-1)}. Bruselas, Bélgica.
    \item \textbf{OpenSees Documentation Team (2024).} \textit{OpenSeesPy Documentation and Command Manual}. UC Berkeley \& PEER Center. \url{https://openseespydoc.readthedocs.io/}
\end{enumerate}

\vspace{0.3cm}
\begin{center}
    {\color{uispurple}\hrule height 1.2pt}\vspace{0.15cm}
    {\footnotesize\color{uisgray} Universidad Industrial de Santander --- Escuela de Ingeniería Civil --- Posgrados --- Bucaramanga, Colombia}
\end{center}

\end{document}
